\rSec1[optional]{Optional objects}

\rSec2[optional.general]{General}

\pnum
Subclause \ref{optional} describes class template \tcode{optional} that represents optional objects. An \defn{optional object} is an object that contains the storage for another object and manages the lifetime of this contained object, if any. The contained object may be initialized after the optional object has been initialized, and may be destroyed before the optional object has been destroyed. The initialization state of the contained object is tracked by the optional object.

\rSec2[optional.syn]{Header \tcode{<optional>} synopsis}

\begin{removedblock}
\pnum
\begin{codeblock}
// mostly freestanding
#include <compare>              // see \ref{compare.syn}

namespace std {
  // \ref{optional.optional}, class template \tcode{optional}
  template<class T>
    class optional;                                             // partially freestanding

  // \ref{optional.optional.ref}, partial specialization of \tcode{optional} for lvalue reference types
  template<class T>
    class optional<T&>;                                         // partially freestanding

  template<class T>
    constexpr bool ranges::enable_view<optional<T>> = true;
  template<class T>
    constexpr auto format_kind<optional<T>> = range_format::disabled;
  template<class T>
    constexpr bool ranges::enable_borrowed_range<optional<T&>> = true;

  template<class T>
    concept @\defexposconceptnc{is-derived-from-optional}@ = requires(const T& t) {   // \expos
      []<class U>(const optional<U>&){ }(t);
    };

  // \ref{optional.nullopt}, no-value state indicator
  struct nullopt_t{@\seebelow@};
  inline constexpr nullopt_t nullopt(@\unspec@);

  // \ref{optional.bad.access}, class \tcode{bad_optional_access}
  class bad_optional_access;

  // \ref{optional.relops}, relational operators
  template<class T, class U>
    constexpr bool operator==(const optional<T>&, const optional<U>&);
  template<class T, class U>
    constexpr bool operator!=(const optional<T>&, const optional<U>&);
  template<class T, class U>
    constexpr bool operator<(const optional<T>&, const optional<U>&);
  template<class T, class U>
    constexpr bool operator>(const optional<T>&, const optional<U>&);
  template<class T, class U>
    constexpr bool operator<=(const optional<T>&, const optional<U>&);
  template<class T, class U>
    constexpr bool operator>=(const optional<T>&, const optional<U>&);
  template<class T, @\libconcept{three_way_comparable_with}@<T> U>
    constexpr compare_three_way_result_t<T, U>
      operator<=>(const optional<T>&, const optional<U>&);

  // \ref{optional.nullops}, comparison with \tcode{nullopt}
  template<class T> constexpr bool operator==(const optional<T>&, nullopt_t) noexcept;
  template<class T>
    constexpr strong_ordering operator<=>(const optional<T>&, nullopt_t) noexcept;

  // \ref{optional.comp.with.t}, comparison with \tcode{T}
  template<class T, class U> constexpr bool operator==(const optional<T>&, const U&);
  template<class T, class U> constexpr bool operator==(const T&, const optional<U>&);
  template<class T, class U> constexpr bool operator!=(const optional<T>&, const U&);
  template<class T, class U> constexpr bool operator!=(const T&, const optional<U>&);
  template<class T, class U> constexpr bool operator<(const optional<T>&, const U&);
  template<class T, class U> constexpr bool operator<(const T&, const optional<U>&);
  template<class T, class U> constexpr bool operator>(const optional<T>&, const U&);
  template<class T, class U> constexpr bool operator>(const T&, const optional<U>&);
  template<class T, class U> constexpr bool operator<=(const optional<T>&, const U&);
  template<class T, class U> constexpr bool operator<=(const T&, const optional<U>&);
  template<class T, class U> constexpr bool operator>=(const optional<T>&, const U&);
  template<class T, class U> constexpr bool operator>=(const T&, const optional<U>&);
  template<class T, class U>
    requires (!@\exposconcept{is-derived-from-optional}@<U>) && @\libconcept{three_way_comparable_with}@<T, U>
    constexpr compare_three_way_result_t<T, U>
      operator<=>(const optional<T>&, const U&);

  // \ref{optional.specalg}, specialized algorithms
  template<class T>
    constexpr void swap(optional<T>&, optional<T>&) noexcept(@\seebelow@);

  template<class T>
    constexpr optional<decay_t<T>> make_optional(T&&);
  template<class T, class... Args>
    constexpr optional<T> make_optional(Args&&... args);
  template<class T, class U, class... Args>
    constexpr optional<T> make_optional(initializer_list<U> il, Args&&... args);

  // \ref{optional.hash}, hash support
  template<class T> struct hash;
  template<class T> struct hash<optional<T>>;
}
\end{codeblock}
\end{removedblock}

\begin{addedblock}
\begin{codeblock}
namespace std {
  // \ref{optional.optional}, class template \tcode{optional}

  template<class T> class @\libglobal{optional}@;

  // \ref{optional.optional.ref}, partial specialization of \tcode{optional} for lvalue reference types

  template<class T> class @\libglobal{optional}@<T&>;

  namespace ranges {
    template<class T> inline constexpr bool ranges::@\libglobal{enable_view}@<optional<T>> = true;

    template<class T>
    inline constexpr bool ranges::@\libglobal{enable_borrowed_range}@<optional<T&>> = true;
  }

  template<class T>
  concept @\exposid{is-derived-from-optional}@ = requires(const T& t) {
    []<class U>(const optional<U>&) {}(t);
  }; // \expos

  // \ref{optional.nullopt}, no-value state indicator

  struct nullopt_t{@\seebelow@};

  inline constexpr nullopt_t @\libglobal{nullopt}@{nullopt_t::Tag::tag};

  // \ref{optional.bad.access}, class \tcode{bad_optional_access}

  class @\libglobal{bad_optional_access}@;

  // \ref{optional.relops}, relational operators

  template<class T, class U>
  constexpr bool @\libglobal{operator==}@(const optional<T>&, const optional<U>&);

  template<class T, class U>
  constexpr bool @\libglobal{operator!=}@(const optional<T>&, const optional<U>&);

  template<class T, class U>
  constexpr bool @\libglobal{operator<}@(const optional<T>&, const optional<U>&);

  template<class T, class U>
  constexpr bool @\libglobal{operator>}@(const optional<T>&, const optional<U>&);

  template<class T, class U>
  constexpr bool @\libglobal{operator<=}@(const optional<T>&, const optional<U>&);

  template<class T, class U>
  constexpr bool @\libglobal{operator>=}@(const optional<T>&, const optional<U>&);

  template<class T, three_way_comparable_with<T> U>
  constexpr compare_three_way_result_t<T, U> @\libglobal{operator<=>}@(const optional<T>&,
                                                         const optional<U>&);

  // \ref{optional.nullops}, comparison with \tcode{nullopt}

  template<class T> constexpr bool @\libglobal{operator==}@(const optional<T>&, nullopt_t) noexcept;

  template<class T>
  constexpr strong_ordering @\libglobal{operator<=>}@(const optional<T>&, nullopt_t) noexcept;

  // \ref{optional.comp.with.t}, comparison with \tcode{T}

  template<class T, class U> constexpr bool @\libglobal{operator==}@(const optional<T>&, const U&);

  template<class T, class U> constexpr bool @\libglobal{operator==}@(const T&, const optional<U>&);

  template<class T, class U> constexpr bool @\libglobal{operator!=}@(const optional<T>&, const U&);

  template<class T, class U> constexpr bool @\libglobal{operator!=}@(const T&, const optional<U>&);

  template<class T, class U> constexpr bool @\libglobal{operator<}@(const optional<T>&, const U&);

  template<class T, class U> constexpr bool @\libglobal{operator<}@(const T&, const optional<U>&);

  template<class T, class U> constexpr bool @\libglobal{operator>}@(const optional<T>&, const U&);

  template<class T, class U> constexpr bool @\libglobal{operator>}@(const T&, const optional<U>&);

  template<class T, class U> constexpr bool @\libglobal{operator<=}@(const optional<T>&, const U&);

  template<class T, class U> constexpr bool @\libglobal{operator<=}@(const T&, const optional<U>&);

  template<class T, class U> constexpr bool @\libglobal{operator>=}@(const optional<T>&, const U&);

  template<class T, class U> constexpr bool @\libglobal{operator>=}@(const T&, const optional<U>&);

  template<class T, class U>
    requires(!@\exposid{is-derived-from-optional}@<U>) && three_way_comparable_with<T, U>
  constexpr compare_three_way_result_t<T, U> @\libglobal{operator<=>}@(const optional<T>&, const U&);

  // \ref{optional.specalg}, specialized algorithms

  template<class T>
  constexpr void @\libglobal{swap}@(optional<T>& x, optional<T>& y) noexcept(noexcept(x.swap(y)));

  template<int = 0, class T>
  constexpr optional<decay_t<T>> @\libglobal{make_optional}@(T&&) noexcept(
      is_nothrow_constructible_v<optional<decay_t<T>>, T>);

  template<class T, class... Args>
  constexpr optional<T> @\libglobal{make_optional}@(Args&&...) noexcept(
      is_nothrow_constructible_v<T, Args...>);

  template<class T, class U, class... Args>
  constexpr optional<T> @\libglobal{make_optional}@(initializer_list<U>, Args&&...) noexcept(
      is_nothrow_constructible_v<T, initializer_list<U>&, Args...>);

  // \ref{optional.hash}, hash support
}
\end{codeblock}
\end{addedblock}

\begin{addedblock}
\rSec2[optional.bad.access]{Class \tcode{bad_optional_access}}

\begin{codeblock}
namespace std {
  class @\libglobal{bad_optional_access}@ : public exception {
  public:
    bad_optional_access() = default;

    const char* what() const noexcept;
  };
}
\end{codeblock}

\pnum
The class \tcode{bad_optional_access} defines the type of objects thrown as exceptions to report the situation where an attempt is made to access the value of an optional object that does not contain a value.

\indexlibrarymember{what}{bad_optional_access}%
\begin{itemdecl}
const char* what() const noexcept;
\end{itemdecl}

\begin{itemdescr}
\pnum
\returns
An \impldef{return value of `bad_optional_access::what`} \ntbs{}, which during constant evaluation is encoded with the ordinary literal encoding\iref{lex.ccon}.
\end{itemdescr}
\end{addedblock}

\rSec2[optional.optional]{Class template \tcode{optional}}

\begin{addedblock}
\begin{codeblock}
namespace std {
  template<class T, class W>
  constexpr bool @\exposid{converts-from-any-cvref}@ = disjunction_v<
      is_constructible<T, W&>, is_convertible<W&, T>, is_constructible<T, W>,
      is_convertible<W, T>, is_constructible<T, const W&>, is_convertible<const W&, T>,
      is_constructible<T, const W>, is_convertible<const W, T>>; // \expos
}
\end{codeblock}
\end{addedblock}

\begin{addedblock}
\pnum
The exposition-only variable template \tcode{converts-from-any-cvref} is used by some constructors for \tcode{optional}.
\end{addedblock}

\rSec3[optional.optional.general]{General}

\begin{codeblock}
namespace std {
  template<class T>
  class @\libglobal{optional}@ {
  public:
    using @\libmember{value_type}{optional}@ = T;
    @\removed{using iterator       = \impdefnc;              // see\textasciitilde{}\ref{optional.iterators}}@
    @\removed{using const_iterator = \impdefnc;              // see\textasciitilde{}\ref{optional.iterators}}@

    @\added{using iterator = \impdefnc; // see\textasciitilde{}\ref{optional.iterators}}@

    @\added{using const_iterator = \impdefnc; // see\textasciitilde{}\ref{optional.iterators}}@

    // \ref{optional.ctor}, constructors
    constexpr optional() noexcept;

    constexpr optional(nullopt_t) noexcept;

    constexpr optional(const optional&);

    constexpr optional(optional&&) noexcept(@\seebelow@);

    template<class... Args> constexpr explicit optional(in_place_t, Args&&...);

    template<class U, class... Args>
    constexpr explicit optional(in_place_t, initializer_list<U>, Args&&...);
    @\removed{template<class U = remove_cv_t<T>>}@
      @\removed{constexpr explicit(\seebelow) optional(U\&\&);}@
    @\removed{template<class U>}@
      @\removed{constexpr explicit(\seebelow) optional(const optional<U>\&);}@
    @\removed{template<class U>}@
      @\removed{constexpr explicit(\seebelow) optional(optional<U>\&\&);}@

    @\added{template<class U = T> constexpr explicit(\seebelow) optional(U\&\&);}@

    @\added{template<class U> constexpr explicit(\seebelow) optional(const optional<U>\&);}@

    @\added{template<class U> constexpr explicit(\seebelow) optional(optional<U>\&\&);}@

    // \ref{optional.dtor}, destructor

    constexpr ~optional();

    // \ref{optional.assign}, assignment
    constexpr optional& operator=(nullopt_t) noexcept;

    constexpr optional& operator=(const optional&);

    constexpr optional& operator=(optional&&) noexcept(@\seebelow@);

    template<class U = @\removed{remove_cv_t<T>>}@@\added{T>}@ constexpr optional& operator=(U&&);

    template<class U> constexpr optional& operator=(const optional<U>&);

    template<class U> constexpr optional& operator=(optional<U>&&);
    template<class... Args> constexpr T& emplace(Args&&...);
    template<class U, class... Args>
    constexpr T& emplace(initializer_list<U>, Args&&...);

    // \ref{optional.swap}, swap
    constexpr void swap(optional&) noexcept(@\seebelow@);

    // \ref{optional.iterators}, iterator support
    constexpr iterator begin() noexcept;

    constexpr const_iterator begin() const noexcept;

    constexpr iterator end() noexcept;

    constexpr const_iterator end() const noexcept;

    // \ref{optional.observe}, observers
    constexpr const T* operator->() const @\removed{noexcept}@;

    constexpr T* operator->()@\removed{noexcept;constexpr const}@@\added{;}@
    @\added{constexpr}@ T& operator*() @\removed{const \& noexcept;constexpr}@@\added{\&;}@

    @\added{constexpr const}@ T& operator*() @\removed{\& noexcept}@@\added{const\&}@;

    constexpr T&& operator*() &&@\removed{noexcept;constexpr const T\&\& operator*() const \&\& noexcept;}@@\added{;}@

    constexpr explicit operator bool() const noexcept;

    constexpr bool has_value() const noexcept;
    @\removed{constexpr const T\& value() const \&;                                 // freestanding-deleted}@
    @\removed{constexpr T\& value() \&;                                             // freestanding-deleted}@
    @\removed{constexpr T\&\& value() \&\&;                                           // freestanding-deleted}@
    @\removed{constexpr const T\&\& value() const \&\&;                               // freestanding-deleted}@
    @\removed{template<class U = remove_cv_t<T>> constexpr T value_or(U\&\&) const \&;}@
    @\removed{template<class U = remove_cv_t<T>> constexpr T value_or(U\&\&) \&\&;}@

    @\added{constexpr T\& value() \&;}@

    @\added{constexpr const T\& value() const\&;}@

    @\added{constexpr T\&\& value() \&\&;}@

    @\added{template<class U = remove_cv_t<T>> constexpr remove_cv_t<T> value_or(U\&\&) const\&;}@
    @\added{template<class U = remove_cv_t<T>> constexpr remove_cv_t<T> value_or(U\&\&) \&\&;}@

    // \ref{optional.monadic}, monadic operations
    @\removed{template<class F> constexpr auto and_then(F\&\& f) \&;}@
    @\removed{template<class F> constexpr auto and_then(F\&\& f) \&\&;}@
    @\removed{template<class F> constexpr auto and_then(F\&\& f) const \&;}@
    @\removed{template<class F> constexpr auto and_then(F\&\& f) const \&\&;}@
    @\removed{template<class F> constexpr auto transform(F\&\& f) \&;}@
    @\removed{template<class F> constexpr auto transform(F\&\& f) \&\&;}@
    @\removed{template<class F> constexpr auto transform(F\&\& f) const \&;}@
    @\removed{template<class F> constexpr auto transform(F\&\& f) const \&\&;}@
    @\removed{template<class F> constexpr optional or_else(F\&\& f) \&\&;}@
    @\removed{template<class F> constexpr optional or_else(F\&\& f) const \&;}@

    @\added{template<class F> constexpr auto and_then(F\&\&) \&;}@

    @\added{template<class F> constexpr auto and_then(F\&\&) \&\&;}@

    @\added{template<class F> constexpr auto and_then(F\&\&) const\&;}@

    @\added{template<class F> constexpr auto and_then(F\&\&) const\&\&;}@

    @\added{template<class F> constexpr auto transform(F\&\&) \&;}@

    @\added{template<class F> constexpr auto transform(F\&\&) \&\&;}@

    @\added{template<class F> constexpr auto transform(F\&\&) const\&;}@

    @\added{template<class F> constexpr auto transform(F\&\&) const\&\&;}@

    @\added{template<class F> constexpr optional or_else(F\&\&) const\&;}@

    @\added{template<class F> constexpr optional or_else(F\&\&) \&\&;}@

    // \ref{optional.mod}, modifiers
    constexpr void reset() noexcept;

  private:
    union {
      @\removed{remove_cv_t<T>}@@\added{T}@ @\exposidnc{val}@; // \expos
    };
  };

  @\removed{template<class T>}@
    @\removed{optional(T) -> optional<T>;}@
}
\end{codeblock}

\pnum
An instance of \tcode{optional<T>} is said to \removed{\defnx{contain a value}{contain a value!\idxcode{optional}}}\added{\defn{contain a value}} when and only when its member \exposid{val} is active\iref{class.union.general}; \exposid{val} is referred to as its \removed{\defnx{contained value}{contained value!\idxcode{optional}}}\added{\defn{contained value}}. An optional object's contained value is nested within\iref{intro.object} the optional object.

\pnum
A type \tcode{X} is a \removed{\defnx{valid contained type}{valid contained type!\idxcode{optional}}}\added{\defn{valid contained type}} for \tcode{optional} if \tcode{X} is an lvalue reference type or a complete non-array object type, and \tcode{remove_cvref_t<X>} is a type other than \tcode{in_place_t} or \tcode{nullopt_t}. If a specialization of \tcode{optional} is instantiated with a type \tcode{T} that is not a valid contained type for \tcode{optional}, the program is ill-formed. If \tcode{T} is an object type, \tcode{T} shall meet the \oldconcept{Destructible} requirements (\tref{cpp17.destructible}).

\begin{addedblock}

\end{addedblock}

\rSec3[optional.ctor]{Constructors}

\begin{removedblock}
\pnum
The exposition-only variable template \exposid{converts-from-any-cvref} is used by some constructors for \tcode{optional}.
\begin{codeblock}
template<class T, class W>
constexpr bool @\exposid{converts-from-any-cvref}@ =  // \expos
  disjunction_v<is_constructible<T, W&>, is_convertible<W&, T>,
                is_constructible<T, W>, is_convertible<W, T>,
                is_constructible<T, const W&>, is_convertible<const W&, T>,
                is_constructible<T, const W>, is_convertible<const W, T>>;
\end{codeblock}
\end{removedblock}

\indexlibraryctor{optional}%
\begin{itemdecl}
constexpr optional() noexcept;
constexpr optional(nullopt_t) noexcept;
\end{itemdecl}

\begin{itemdescr}
\pnum
\ensures
\tcode{*this} does not contain a value.

\pnum
\remarks
No contained value is initialized. For every object type \tcode{T} these constructors are constexpr constructors\iref{dcl.constexpr}.
\end{itemdescr}

\indexlibraryctor{optional}%
\begin{itemdecl}
constexpr optional(const optional& rhs);
\end{itemdecl}

\begin{itemdescr}
\begin{addedblock}

\end{addedblock}

\pnum
\effects
If \tcode{rhs} contains a value, direct-non-list-initializes \exposid{val} with \tcode{rhs.\exposid{val}}.

\pnum
\ensures
\tcode{rhs.has_value() == this->has_value()}.

\pnum
\throws
Any exception thrown by the selected constructor of \tcode{T}.

\pnum
\remarks
This constructor is defined as deleted unless \tcode{is_copy_constructible_v<T>} is \tcode{true}. If \tcode{is_trivially_copy_constructible_v<T>} is \tcode{true}, this constructor is trivial.
\end{itemdescr}

\indexlibraryctor{optional}%
\begin{itemdecl}
constexpr optional(optional&& rhs) noexcept(@\seebelow@);
\end{itemdecl}

\begin{itemdescr}
\pnum
\constraints
\tcode{is_move_constructible_v<T>} is \tcode{true}.

\pnum
\effects
If \tcode{rhs} contains a value, direct-non-list-initializes \exposid{val} with \tcode{std::move(rhs.\exposid{val})}. \tcode{rhs.has_value()} is unchanged.

\pnum
\ensures
\tcode{rhs.has_value() == this->has_value()}.

\pnum
\throws
Any exception thrown by the selected constructor of \tcode{T}.

\pnum
\remarks
The exception specification is equivalent to \tcode{is_nothrow_move_constructible_v<T>}. If \tcode{is_trivially_move_constructible_v<T>} is \tcode{true}, this constructor is trivial.
\end{itemdescr}

\indexlibraryctor{optional}%
\begin{itemdecl}
template<class... Args>
  constexpr explicit optional(in_place_t, Args&&... args);
\end{itemdecl}

\begin{itemdescr}
\pnum
\constraints
\tcode{is_constructible_v<T, Args...>} is \tcode{true}.

\pnum
\effects
Direct-non-list-initializes \exposid{val} with \tcode{std::forward<Args>(args)...}.

\pnum
\ensures
\tcode{*this} contains a value.

\pnum
\throws
Any exception thrown by the selected constructor of \tcode{T}.

\pnum
\remarks
If \tcode{T}'s constructor selected for the initialization is a constexpr constructor, this constructor is a constexpr constructor.
\end{itemdescr}

\indexlibraryctor{optional}%
\begin{itemdecl}
template<class U, class... Args>
  constexpr explicit optional(in_place_t, initializer_list<U> il, Args&&... args);
\end{itemdecl}

\begin{itemdescr}
\pnum
\constraints
\tcode{\removed{is_constructible_v<T, initializer_list<U>\&, Args...>}}\tcode{\added{is_constructible_v<T, initializer_list<U>\&, Args\&\&...>}} is \tcode{true}.

\pnum
\effects
Direct-non-list-initializes \exposid{val} with \tcode{il, std::forward<Args>(args)...}.

\pnum
\ensures
\tcode{*this} contains a value.

\pnum
\throws
Any exception thrown by the selected constructor of \tcode{T}.

\pnum
\remarks
If \tcode{T}'s constructor selected for the initialization is a constexpr constructor, this constructor is a constexpr constructor.
\end{itemdescr}

\indexlibraryctor{optional}%
\begin{itemdecl}
template<class U = @\removed{remove_cv_t<T>>}@@\added{T>}@
  constexpr explicit(@\seebelow@) optional(U&& v);
\end{itemdecl}

\begin{itemdescr}
\pnum
\constraints
\begin{itemize}
\item \tcode{is_constructible_v<T, U>} is \tcode{true},
\item \tcode{is_same_v<remove_cvref_t<U>, in_place_t>} is \tcode{false}, \added{and}
\item \tcode{is_same_v<remove_cvref_t<U>, optional>} is \tcode{false}\removed{, and}\added{.}
\item \removed{if }\tcode{\removed{T}}\removed{ is \cv{} }\tcode{\removed{bool}}\removed{, }\tcode{\removed{remove_cvref_t<U>}}\removed{ is not a specialization of }\tcode{\removed{optional}}\removed{.}
\end{itemize}

\pnum
\effects
Direct-non-list-initializes \exposid{val} with \tcode{std::forward<U>(v)}.

\pnum
\ensures
\tcode{*this} contains a value.

\pnum
\throws
Any exception thrown by the selected constructor of \tcode{T}.

\pnum
\remarks
If \tcode{T}'s selected constructor is a constexpr constructor, this constructor is a constexpr constructor. The expression inside \removed{\keyword{explicit}}\tcode{\added{explicit}} is equivalent to:
\begin{codeblock}
!is_convertible_v<U, T>
\end{codeblock}
\end{itemdescr}

\indexlibraryctor{optional}%
\begin{itemdecl}
template<class U>
  constexpr explicit(@\seebelow@) optional(const optional<U>& rhs);
\end{itemdecl}

\begin{itemdescr}
\pnum
\constraints
\begin{itemize}
\item \tcode{\added{is_same_v<T, U>}}\added{ is }\tcode{\added{false}}\added{,}
\item \tcode{is_constructible_v<T, const U\&>} is \tcode{true}, and
\item \removed{if }\tcode{\removed{T}}\removed{ is not \cv{} }\tcode{\removed{bool}}\removed{, }\tcode{\removed{\exposid{converts-from-any-cvref}<T, optional<U>>}}\removed{ is }\tcode{\removed{false}}\removed{.}
\item \tcode{\added{converts-from-any-cvref<T, optional<U>>}}\added{ is }\tcode{\added{false}}\added{.}
\end{itemize}

\pnum
\effects
If \tcode{rhs} contains a value, direct-non-list-initializes \exposid{val} with \tcode{rhs.operator*()}.

\pnum
\ensures
\tcode{rhs.has_value() == this->has_value()}.

\pnum
\throws
Any exception thrown by the selected constructor of \tcode{T}.

\pnum
\remarks
The expression inside \removed{\keyword{explicit}}\tcode{\added{explicit}} is equivalent to:
\begin{codeblock}
!is_convertible_v<@\removed{const U\&}@@\added{U}@, T>
\end{codeblock}
\end{itemdescr}

\indexlibraryctor{optional}%
\begin{itemdecl}
template<class U>
  constexpr explicit(@\seebelow@) optional(optional<U>&& rhs);
\end{itemdecl}

\begin{itemdescr}
\pnum
\constraints
\begin{itemize}
\item \tcode{\added{is_same_v<T, U>}}\added{ is }\tcode{\added{false}}\added{,}
\item \tcode{is_constructible_v<T, U>} is \tcode{true}, and
\item \removed{if }\tcode{\removed{T}}\removed{ is not \cv{} }\tcode{\removed{bool}}\removed{, }\tcode{\removed{\exposid{converts-from-any-cvref}<T, optional<U>>}}\removed{ is }\tcode{\removed{false}}\removed{.}
\item \tcode{\added{converts-from-any-cvref<T, optional<U>>}}\added{ is }\tcode{\added{false}}\added{.}
\end{itemize}

\pnum
\effects
If \tcode{rhs} contains a value, direct-non-list-initializes \exposid{val} with \tcode{std::move(rhs).operator*()}. \tcode{rhs.has_value()} is unchanged.

\pnum
\ensures
\tcode{rhs.has_value() == this->has_value()}.

\pnum
\throws
Any exception thrown by the selected constructor of \tcode{T}.

\pnum
\remarks
The expression inside \removed{\keyword{explicit}}\tcode{\added{explicit}} is equivalent to:
\begin{codeblock}
!is_convertible_v<U, T>
\end{codeblock}
\end{itemdescr}

\rSec3[optional.dtor]{Destructor}

\indexlibrarydtor{optional}%
\begin{itemdecl}
constexpr ~optional();
\end{itemdecl}

\begin{itemdescr}
\pnum
\effects
If \tcode{is_trivially_destructible_v<T>} is \tcode{false} and \tcode{*this} contains a value, calls \tcode{\exposid{val}.T::\textasciitilde{}T()}.

\pnum
\remarks
If \tcode{is_trivially_destructible_v<T>} is \tcode{true}, then this destructor is trivial.
\end{itemdescr}

\rSec3[optional.assign]{Assignment}

\indexlibrarymember{operator=}{optional}%
\begin{itemdecl}
constexpr optional@\removed{<T>}@& operator=(nullopt_t) noexcept;
\end{itemdecl}

\begin{itemdescr}
\pnum
\effects
If \tcode{*this} contains a value, calls \tcode{\exposid{val}.T::\textasciitilde{}T()} to destroy the contained value; otherwise no effect.

\pnum
\ensures
\tcode{*this} does not contain a value.

\pnum
\returns
\tcode{*this}.
\end{itemdescr}

\indexlibrarymember{operator=}{optional}%
\begin{itemdecl}
constexpr optional@\removed{<T>}@& operator=(const optional& rhs);
\end{itemdecl}

\begin{itemdescr}
\begin{addedblock}

\end{addedblock}

\pnum
\effects
See \tref{optional.assign.copy}.
\begin{lib2dtab2}{\tcode{optional::operator=(const optional\&)} effects}{optional.assign.copy}
{\tcode{*this} contains a value}
{\tcode{*this} does not contain a value}

\rowhdr{\tcode{rhs} contains a value} &
assigns \tcode{rhs.\exposid{val}} to \exposid{val} &
direct-non-list-initializes \exposid{val} with \tcode{rhs.\exposid{val}} \\
\rowsep

\rowhdr{\tcode{rhs} does not contain a value} &
destroys the contained value by calling \tcode{\exposid{val}.T::\textasciitilde{}T()} &
no effect \\
\end{lib2dtab2}

\pnum
\ensures
\tcode{rhs.has_value() == this->has_value()}.

\pnum
\returns
\tcode{*this}.

\pnum
\remarks
If any exception is thrown, the result of the expression \tcode{this->has_value()} remains unchanged. If an exception is thrown during the call to \tcode{T}'s copy constructor, no effect. If an exception is thrown during the call to \tcode{T}'s copy assignment, the state of its contained value is as defined by the exception safety guarantee of \tcode{T}'s copy assignment. This operator is defined as deleted unless \tcode{is_copy_constructible_v<T>} is \tcode{true} and \tcode{is_copy_assignable_v<T>} is \tcode{true}. If \tcode{is_trivially_copy_constructible_v<T> \&\&} \tcode{is_trivially_copy_assignable_v<T> \&\&} \tcode{is_trivially_destructible_v<T>} is \tcode{true}, this assignment operator is trivial.
\end{itemdescr}

\indexlibrarymember{operator=}{optional}%
\begin{itemdecl}
constexpr optional& operator=(optional&& rhs) noexcept(@\seebelow@);
\end{itemdecl}

\begin{itemdescr}
\pnum
\constraints
\tcode{is_move_constructible_v<T>} is \tcode{true} and \tcode{is_move_assignable_v<T>} is \tcode{true}.

\pnum
\effects
See \tref{optional.assign.move}. The result of the expression \tcode{rhs.has_value()} remains unchanged.
\begin{lib2dtab2}{\tcode{optional::operator=(optional\&\&)} effects}{optional.assign.move}
{\tcode{*this} contains a value}
{\tcode{*this} does not contain a value}

\rowhdr{\tcode{rhs} contains a value} &
assigns \tcode{std::move(rhs.\exposid{val})} to \exposid{val} &
direct-non-list-initializes \exposid{val} with \tcode{std::move(rhs.\exposid{val})} \\
\rowsep

\rowhdr{\tcode{rhs} does not contain a value} &
destroys the contained value by calling \tcode{\exposid{val}.T::\textasciitilde{}T()} &
no effect \\
\end{lib2dtab2}

\pnum
\ensures
\tcode{rhs.has_value() == this->has_value()}.

\pnum
\returns
\tcode{*this}.

\pnum
\remarks
The exception specification is equivalent to:
\begin{codeblock}
@\removed{is_nothrow_move_assignable_v<T> \&\&}@ is_nothrow_move_constructible_v<T>
\end{codeblock}

\pnum
If any exception is thrown, the result of the expression \tcode{this->has_value()} remains unchanged. If an exception is thrown during the call to \tcode{T}'s move constructor, the state of \tcode{rhs.\exposid{val}} is determined by the exception safety guarantee of \tcode{T}'s move constructor. If an exception is thrown during the call to \tcode{T}'s move assignment, the states of \exposid{val} and \tcode{rhs.\exposid{val}} are determined by the exception safety guarantee of \tcode{T}'s move assignment. If \tcode{is_trivially_move_constructible_v<T> \&\&} \tcode{is_trivially_move_assignable_v<T> \&\&} \tcode{is_trivially_destructible_v<T>} is \tcode{true}, this assignment operator is trivial.
\end{itemdescr}

\indexlibrarymember{operator=}{optional}%
\begin{itemdecl}
template<class U = @\removed{remove_cv_t<T>>}@@\added{T>}@
  constexpr optional& operator=(U&& v);
\end{itemdecl}

\begin{itemdescr}
\pnum
\constraints
\begin{itemize}
\item \tcode{is_same_v<remove_cvref_t<U>, optional>} is \tcode{false},
\item \tcode{conjunction_v<is_scalar<T>, is_same<T, decay_t<U>>>} is \tcode{false},
\item \tcode{is_constructible_v<T, U>} is \tcode{true}, and
\item \tcode{is_assignable_v<T\&, U>} is \tcode{true}.
\end{itemize}

\pnum
\effects
If \tcode{*this} contains a value, assigns \tcode{std::forward<U>(v)} to \exposid{val}; otherwise direct-non-list-initializes \exposid{val} with \tcode{std::forward<U>(v)}.

\pnum
\ensures
\tcode{*this} contains a value.

\pnum
\returns
\tcode{*this}.

\pnum
\remarks
If any exception is thrown, the result of the expression \tcode{this->has_value()} remains unchanged. If an exception is thrown during the call to \tcode{T}'s constructor, the state of \tcode{v} is determined by the exception safety guarantee of \tcode{T}'s constructor. If an exception is thrown during the call to \tcode{T}'s assignment, the states of \exposid{val} and \tcode{v} are determined by the exception safety guarantee of \tcode{T}'s assignment.
\end{itemdescr}

\indexlibrarymember{operator=}{optional}%
\begin{itemdecl}
template<class U>
  constexpr optional@\removed{<T>}@& operator=(const optional<U>& rhs);
\end{itemdecl}

\begin{itemdescr}
\pnum
\constraints
\begin{itemize}
\item \tcode{is_constructible_v<T, const U\&>} is \tcode{true},
\item \tcode{is_assignable_v<T\&, const U\&>} is \tcode{true},
\item \tcode{\removed{\exposid{converts-from-any-cvref}<T, optional<U>>}}\tcode{\added{converts-from-any-cvref<T, optional<U>>}} is \tcode{false},
\item \tcode{is_assignable_v<T\&, optional<U>\&>} is \tcode{false},
\item \tcode{is_assignable_v<T\&, optional<U>\&\&>} is \tcode{false},
\item \tcode{is_assignable_v<T\&, const optional<U>\&>} is \tcode{false}, and
\item \tcode{is_assignable_v<T\&, const optional<U>\&\&>} is \tcode{false}.
\end{itemize}

\pnum
\effects
See \tref{optional.assign.copy.templ}.
\begin{lib2dtab2}{\tcode{optional::operator=(const optional<U>\&)} effects}{optional.assign.copy.templ}
{\tcode{*this} contains a value}
{\tcode{*this} does not contain a value}

\rowhdr{\tcode{rhs} contains a value} &
assigns \tcode{rhs.operator*()} to \exposid{val} &
direct-non-list-initializes \exposid{val} with \tcode{rhs.operator*()} \\
\rowsep

\rowhdr{\tcode{rhs} does not contain a value} &
destroys the contained value by calling \tcode{\exposid{val}.T::\textasciitilde{}T()} &
no effect \\
\end{lib2dtab2}

\pnum
\ensures
\tcode{rhs.has_value() == this->has_value()}.

\pnum
\returns
\tcode{*this}.

\pnum
\remarks
If any exception is thrown, the result of the expression \tcode{this->has_value()} remains unchanged. If an exception is thrown during the call to \tcode{T}'s constructor, the state of \tcode{rhs.\exposid{val}} is determined by the exception safety guarantee of \tcode{T}'s constructor. If an exception is thrown during the call to \tcode{T}'s assignment, the states of \exposid{val} and \tcode{rhs.\exposid{val}} are determined by the exception safety guarantee of \tcode{T}'s assignment.
\end{itemdescr}

\indexlibrarymember{operator=}{optional}%
\begin{itemdecl}
template<class U>
  constexpr optional@\removed{<T>}@& operator=(optional<U>&& rhs);
\end{itemdecl}

\begin{itemdescr}
\pnum
\constraints
\begin{itemize}
\item \tcode{is_constructible_v<T, U>} is \tcode{true},
\item \tcode{is_assignable_v<T\&, U>} is \tcode{true},
\item \tcode{\removed{\exposid{converts-from-any-cvref}<T, optional<U>>}}\tcode{\added{converts-from-any-cvref<T, optional<U>>}} is \tcode{false},
\item \tcode{is_assignable_v<T\&, optional<U>\&>} is \tcode{false},
\item \tcode{is_assignable_v<T\&, optional<U>\&\&>} is \tcode{false},
\item \tcode{is_assignable_v<T\&, const optional<U>\&>} is \tcode{false}, and
\item \tcode{is_assignable_v<T\&, const optional<U>\&\&>} is \tcode{false}.
\end{itemize}

\pnum
\effects
See \tref{optional.assign.move.templ}. The result of the expression \tcode{rhs.has_value()} remains unchanged.
\begin{lib2dtab2}{\tcode{optional::operator=(optional<U>\&\&)} effects}{optional.assign.move.templ}
{\tcode{*this} contains a value}
{\tcode{*this} does not contain a value}

\rowhdr{\tcode{rhs} contains a value} &
assigns \tcode{std::move(rhs).operator*()} to \exposid{val} &
direct-non-list-initializes \exposid{val} with \tcode{std::move(rhs).operator*()} \\
\rowsep

\rowhdr{\tcode{rhs} does not contain a value} &
destroys the contained value by calling \tcode{\exposid{val}.T::\textasciitilde{}T()} &
no effect \\
\end{lib2dtab2}

\pnum
\ensures
\tcode{rhs.has_value() == this->has_value()}.

\pnum
\returns
\tcode{*this}.

\pnum
\remarks
If any exception is thrown, the result of the expression \tcode{this->has_value()} remains unchanged. If an exception is thrown during the call to \tcode{T}'s constructor, the state of \tcode{rhs.\exposid{val}} is determined by the exception safety guarantee of \tcode{T}'s constructor. If an exception is thrown during the call to \tcode{T}'s assignment, the states of \exposid{val} and \tcode{rhs.\exposid{val}} are determined by the exception safety guarantee of \tcode{T}'s assignment.
\end{itemdescr}

\indexlibrarymember{emplace}{optional}%
\begin{itemdecl}
template<class... Args>
  constexpr T& emplace(Args&&... args);
\end{itemdecl}

\begin{itemdescr}
\begin{removedblock}
\pnum
\constraints
\tcode{is_constructible_v<T, Args...>} is \tcode{true}.
\end{removedblock}

\begin{addedblock}
\pnum
\mandates
\tcode{is_constructible_v<T, Args\&\&...>} is \tcode{true}.
\end{addedblock}

\pnum
\effects
Calls \tcode{*this = nullopt}. Then direct-non-list-initializes \exposid{val} with \tcode{std::forward<Args>(args)...}.

\pnum
\ensures
\tcode{*this} contains a value.

\pnum
\returns
\exposid{val}.

\pnum
\throws
Any exception thrown by the selected constructor of \tcode{T}.

\pnum
\remarks
If an exception is thrown during the call to \tcode{T}'s constructor, \tcode{*this} does not contain a value, and the previous \exposid{val} (if any) has been destroyed.
\end{itemdescr}

\indexlibrarymember{emplace}{optional}%
\begin{itemdecl}
template<class U, class... Args>
  constexpr T& emplace(initializer_list<U> il, Args&&... args);
\end{itemdecl}

\begin{itemdescr}
\begin{removedblock}
\pnum
\constraints
\tcode{is_constructible_v<T, initializer_list<U>\&, Args...>} is \tcode{true}.
\end{removedblock}

\begin{addedblock}
\pnum
\mandates
\tcode{is_constructible_v<T, initializer_list<U>\&, Args\&\&...>} is \tcode{true}.
\end{addedblock}

\pnum
\effects
Calls \tcode{*this = nullopt}. Then direct-non-list-initializes \exposid{val} with \tcode{il, std::forward<Args>(args)...}.

\pnum
\ensures
\tcode{*this} contains a value.

\pnum
\returns
\exposid{val}.

\pnum
\throws
Any exception thrown by the selected constructor of \tcode{T}.

\pnum
\remarks
If an exception is thrown during the call to \tcode{T}'s constructor, \tcode{*this} does not contain a value, and the previous \exposid{val} (if any) has been destroyed.
\end{itemdescr}

\rSec3[optional.swap]{Swap}

\indexlibrarymember{swap}{optional}%
\begin{itemdecl}
constexpr void swap(optional& rhs) noexcept(@\seebelow@);
\end{itemdecl}

\begin{itemdescr}
\pnum
\mandates
\tcode{is_move_constructible_v<T>} is \tcode{true}.

\pnum
\expects
\tcode{T} meets the \oldconcept{Swappable} requirements\iref{swappable.requirements}.

\pnum
\effects
See \tref{optional.swap}.
\begin{lib2dtab2}{\tcode{optional::swap(optional\&)} effects}{optional.swap}
{\tcode{*this} contains a value}
{\tcode{*this} does not contain a value}

\rowhdr{\tcode{rhs} contains a value} &
calls \tcode{swap(\exposid{val}, rhs.\exposid{val})} &
direct-non-list-initializes \exposid{val} with \tcode{std::move(rhs.\exposid{val})}, followed by \tcode{rhs.\exposid{val}.T::\textasciitilde{}T()}; postcondition is that \tcode{*this} contains a value and \tcode{rhs} does not contain a value \\
\rowsep

\rowhdr{\tcode{rhs} does not contain a value} &
direct-non-list-initializes \tcode{rhs.\exposid{val}} with \tcode{std::move(\exposid{val})}, followed by \tcode{\exposid{val}.T::\textasciitilde{}T()}; postcondition is that \tcode{*this} does not contain a value and \tcode{rhs} contains a value &
no effect \\
\end{lib2dtab2}

\pnum
\throws
Any exceptions thrown by the operations in the relevant part of \tref{optional.swap}.

\pnum
\remarks
The exception specification is equivalent to:
\begin{codeblock}
is_nothrow_move_constructible_v<T> && is_nothrow_swappable_v<T>
\end{codeblock}

\pnum
If any exception is thrown, the results of the expressions \tcode{this->has_value()} and \tcode{rhs.has_value()} remain unchanged. If an exception is thrown during the call to function \tcode{swap}, the states of \exposid{val} and \tcode{rhs.\exposid{val}} are determined by the exception safety guarantee of \tcode{swap} for lvalues of \tcode{T}. If an exception is thrown during the call to \tcode{T}'s move constructor, the states of \exposid{val} and \tcode{rhs.\exposid{val}} are determined by the exception safety guarantee of \tcode{T}'s move constructor.
\end{itemdescr}

\rSec3[optional.iterators]{Iterator support}

\indexlibrarymember{iterator}{optional}%
\indexlibrarymember{const_iterator}{optional}%
\begin{itemdecl}
using iterator = @\impdef@;
using const_iterator = @\impdef@;
\end{itemdecl}

\begin{itemdescr}
\pnum
These types model \libconcept{contiguous_iterator}\iref{iterator.concept.contiguous}, meet the \oldconcept{RandomAccessIterator} requirements\iref{random.access.iterators}, and meet the requirements for constexpr iterators\iref{iterator.requirements.general}, with value type \tcode{remove_cv_t<T>}. The reference type is \tcode{T\&} for \tcode{iterator} and \tcode{const T\&} for \tcode{const_iterator}.

\pnum
All requirements on container iterators\iref{container.reqmts} apply to \tcode{optional::iterator} and \tcode{optional::const_iterator} as well.

\pnum
Any operation that initializes or destroys the contained value of an optional object invalidates all iterators into that object.
\end{itemdescr}

\indexlibrarymember{begin}{optional}%
\begin{itemdecl}
constexpr iterator begin() noexcept;
constexpr const_iterator begin() const noexcept;
\end{itemdecl}

\begin{itemdescr}
\pnum
\returns
If \tcode{has_value()} is \tcode{true}, an iterator referring to the contained value. Otherwise, a past-the-end iterator value.
\end{itemdescr}

\indexlibrarymember{end}{optional}%
\begin{itemdecl}
constexpr iterator end() noexcept;
constexpr const_iterator end() const noexcept;
\end{itemdecl}

\begin{itemdescr}
\pnum
\returns
\tcode{begin() + has_value()}.
\end{itemdescr}

\rSec3[optional.observe]{Observers}

\indexlibrarymember{operator->}{optional}%
\begin{itemdecl}
constexpr const T* operator->() const @\removed{noexcept}@;
constexpr T* operator->() @\removed{noexcept}@;
\end{itemdecl}

\begin{itemdescr}
\pnum
\hardexpects
\tcode{has_value()} is \tcode{true}.

\pnum
\returns
\tcode{addressof(\exposid{val})}.

\pnum
\remarks
These functions are constexpr functions.
\end{itemdescr}

\indexlibrarymember{operator*}{optional}%
\begin{itemdecl}
constexpr @\removed{const}@ T& operator*() @\removed{const \& noexcept}@@\added{\&}@;
constexpr @\added{const}@ T& operator*() @\removed{\& noexcept}@@\added{const\&}@;
\end{itemdecl}

\begin{itemdescr}
\pnum
\hardexpects
\tcode{has_value()} is \tcode{true}.

\pnum
\returns
\exposid{val}.

\pnum
\remarks
These functions are constexpr functions.
\end{itemdescr}

\indexlibrarymember{operator*}{optional}%
\begin{itemdecl}
constexpr T&& operator*() && @\removed{noexcept}@;
@\removed{constexpr const T\&\& operator*() const \&\& noexcept;}@
\end{itemdecl}

\begin{itemdescr}
\pnum
\hardexpects
\tcode{has_value()} is \tcode{true}.

\begin{removedblock}
\pnum
\effects
Equivalent to: \tcode{return std::move(\exposid{val});}
\end{removedblock}

\begin{addedblock}
\pnum
Equivalent to:
\begin{codeblock}
return std::move(@\exposid{val}@);
\end{codeblock}
\end{addedblock}
\end{itemdescr}

\indexlibrarymember{operator bool}{optional}%
\begin{itemdecl}
constexpr explicit operator bool() const noexcept;
\end{itemdecl}

\begin{itemdescr}
\pnum
\returns
\tcode{true} if and only if \tcode{*this} contains a value.

\pnum
\remarks
This function is a constexpr function.
\end{itemdescr}

\indexlibrarymember{has_value}{optional}%
\begin{itemdecl}
constexpr bool has_value() const noexcept;
\end{itemdecl}

\begin{itemdescr}
\pnum
\returns
\tcode{true} if and only if \tcode{*this} contains a value.

\pnum
\remarks
This function is a constexpr function.
\end{itemdescr}

\indexlibrarymember{value}{optional}%
\begin{itemdecl}
constexpr @\removed{const}@ T& value() @\removed{const}@ &;
constexpr @\added{const}@ T& value() @\added{const}@&;
\end{itemdecl}

\begin{itemdescr}
\pnum
\effects
Equivalent to:
\begin{codeblock}
return has_value() ? @\exposid{val}@ : throw bad_optional_access();
\end{codeblock}
\end{itemdescr}

\indexlibrarymember{value}{optional}%
\begin{itemdecl}
constexpr T&& value() &&;
@\removed{constexpr const T\&\& value() const \&\&;}@
\end{itemdecl}

\begin{itemdescr}
\pnum
\effects
Equivalent to:
\begin{codeblock}
return has_value() ? std::move(@\exposid{val}@) : throw bad_optional_access();
\end{codeblock}
\end{itemdescr}

\indexlibrarymember{value_or}{optional}%
\begin{itemdecl}
template<class U = remove_cv_t<T>>
  constexpr @\removed{T}@@\added{remove_cv_t<T>}@ value_or(U&& v) const&;
\end{itemdecl}

\begin{itemdescr}
\pnum
\mandates
\tcode{\removed{is_copy_constructible_v<T> \&\& is_convertible_v<U\&\&, T>}}\tcode{\added{is_convertible_v<const T\&, X>}} is \tcode{true} \added{and }\tcode{\added{is_convertible_v<U, X>}}\added{ is }\tcode{\added{true}}.

\pnum
\effects
Equivalent to:
\begin{codeblock}
@\removed{return has_value() ? \exposid{val} : static_cast<T>(std::forward<U>(v));}@
@\added{using X = remove_cv_t<T>;}@

@\added{if (has_value())}@
  @\added{return \exposid{val};}@
@\added{return std::forward<U>(v);}@
\end{codeblock}
\end{itemdescr}

\indexlibrarymember{value_or}{optional}%
\begin{itemdecl}
template<class U = remove_cv_t<T>>
  constexpr @\removed{T}@@\added{remove_cv_t<T>}@ value_or(U&& v) &&;
\end{itemdecl}

\begin{itemdescr}
\pnum
\mandates
\tcode{\removed{is_move_constructible_v<T> \&\& is_convertible_v<U\&\&, T>}}\tcode{\added{is_convertible_v<T, X>}} is \tcode{true} \added{and }\tcode{\added{is_convertible_v<U, X>}}\added{ is }\tcode{\added{true}}.

\pnum
\effects
Equivalent to:
\begin{codeblock}
@\removed{return has_value() ? std::move(\exposid{val}) : static_cast<T>(std::forward<U>(v));}@
@\added{using X = remove_cv_t<T>;}@

@\added{if (has_value()) \{}@
  @\added{return std::move(\exposid{val});}@
@\added{\}}@
@\added{return std::forward<U>(v);}@
\end{codeblock}
\end{itemdescr}

\rSec3[optional.monadic]{Monadic operations}

\indexlibrarymember{and_then}{optional}%
\begin{itemdecl}
template<class F>
  constexpr auto and_then(F&& f) &;
template<class F>
  constexpr auto and_then(F&& f) const&;
\end{itemdecl}

\begin{itemdescr}
\pnum
Let \tcode{U} be \tcode{invoke_result_t<F, decltype((\exposid{val}))>}.

\pnum
\mandates
\tcode{remove_cvref_t<U>} is a specialization of \tcode{optional}.

\pnum
\effects
Equivalent to:
\begin{codeblock}
if (*this) {
  return invoke(std::forward<F>(f), @\exposid{val}@);
} else {
  return remove_cvref_t<U>();
}
\end{codeblock}
\end{itemdescr}

\indexlibrarymember{and_then}{optional}%
\begin{itemdecl}
template<class F>
  constexpr auto and_then(F&& f) &&;
template<class F>
  constexpr auto and_then(F&& f) const&&;
\end{itemdecl}

\begin{itemdescr}
\pnum
Let \tcode{U} be \tcode{invoke_result_t<F, decltype(std::move(\exposid{val}))>}.

\pnum
\mandates
\tcode{remove_cvref_t<U>} is a specialization of \tcode{optional}.

\pnum
\effects
Equivalent to:
\begin{codeblock}
if (*this) {
  return invoke(std::forward<F>(f), std::move(@\exposid{val}@));
} else {
  return remove_cvref_t<U>();
}
\end{codeblock}
\end{itemdescr}

\indexlibrarymember{transform}{optional}%
\begin{itemdecl}
template<class F>
  constexpr auto transform(F&& f) &;
template<class F>
  constexpr auto transform(F&& f) const&;
\end{itemdecl}

\begin{itemdescr}
\pnum
Let \tcode{U} be \tcode{\removed{remove_cv_t<invoke_result_t<F, decltype((\exposid{val}))>>}}\tcode{\added{invoke_result_t<F, decltype((\exposid{val}))>}}.

\pnum
\mandates
\tcode{U} is a valid contained type for \tcode{optional}. The declaration
\begin{codeblock}
U u(invoke(std::forward<F>(f), @\exposid{val}@));
\end{codeblock}
is well-formed for some invented variable \tcode{u}.
\begin{note}
There is no requirement that \tcode{U} is movable\iref{dcl.init.general}.
\end{note}

\pnum
\returns
If \tcode{*this} contains a value, an \tcode{optional<U>} object whose contained value is direct-non-list-initialized with \tcode{invoke(std::forward<F>(f), \exposid{val})}; otherwise, \tcode{optional<U>()}.
\end{itemdescr}

\indexlibrarymember{transform}{optional}%
\begin{itemdecl}
template<class F>
  constexpr auto transform(F&& f) &&;
template<class F>
  constexpr auto transform(F&& f) const&&;
\end{itemdecl}

\begin{itemdescr}
\pnum
Let \tcode{U} be \tcode{\removed{remove_cv_t<invoke_result_t<F, decltype(std::move(\exposid{val}))>>}}\tcode{\added{invoke_result_t<F, decltype(std::move(\exposid{val}))>}}.

\pnum
\mandates
\tcode{U} is a valid contained type for \tcode{optional}. The declaration
\begin{codeblock}
U u(invoke(std::forward<F>(f), std::move(@\exposid{val}@)));
\end{codeblock}
is well-formed for some invented variable \tcode{u}.
\begin{note}
There is no requirement that \tcode{U} is movable\iref{dcl.init.general}.
\end{note}

\pnum
\returns
If \tcode{*this} contains a value, an \tcode{optional<U>} object whose contained value is direct-non-list-initialized with \tcode{invoke(std::forward<F>(f), std::move(\exposid{val}))}; otherwise, \tcode{optional<U>()}.
\end{itemdescr}

\indexlibrarymember{or_else}{optional}%
\begin{itemdecl}
template<class F>
  constexpr optional or_else(F&& f) const&;
\end{itemdecl}

\begin{itemdescr}
\pnum
\constraints
\tcode{F} models \libconcept{invocable} and \tcode{T} models \libconcept{copy_constructible}.

\pnum
\mandates
\tcode{is_same_v<remove_cvref_t<invoke_result_t<F>>, optional>} is \tcode{true}.

\pnum
\effects
Equivalent to:
\begin{codeblock}
@\removed{if (*this) \{}@
  return @\removed{*this}@@\added{\exposid{val}}@;
@\removed{\} else \{}@
  return std::forward<F>(f)();
@\removed{\}}@
@\added{if (has_value())}@

\end{codeblock}
\end{itemdescr}

\indexlibrarymember{or_else}{optional}%
\begin{itemdecl}
template<class F>
  constexpr optional or_else(F&& f) &&;
\end{itemdecl}

\begin{itemdescr}
\pnum
\constraints
\tcode{F} models \libconcept{invocable} and \tcode{T} models \libconcept{move_constructible}.

\pnum
\mandates
\tcode{is_same_v<remove_cvref_t<invoke_result_t<F>>, optional>} is \tcode{true}.

\pnum
\effects
Equivalent to:
\begin{codeblock}
@\removed{if (*this) \{}@
  return std::move(@\removed{*this}@@\added{\exposid{val}}@);
@\removed{\} else \{}@
  return std::forward<F>(f)();
@\removed{\}}@
@\added{if (has_value())}@

\end{codeblock}
\end{itemdescr}

\rSec3[optional.mod]{Modifiers}

\indexlibrarymember{reset}{optional}%
\begin{itemdecl}
constexpr void reset() noexcept;
\end{itemdecl}

\begin{itemdescr}
\pnum
\effects
If \tcode{*this} contains a value, calls \tcode{\exposid{val}.T::\textasciitilde{}T()} to destroy the contained value; otherwise no effect.

\pnum
\ensures
\tcode{*this} does not contain a value.
\end{itemdescr}

\rSec2[optional.optional.ref]{Partial specialization of \tcode{optional} for reference types}

\rSec3[optional.optional.ref.general]{General}

\begin{codeblock}
namespace std {
  template<class T>
  class @\libglobal{optional}@<T&> {
  public:
    using @\libmember{value_type}{optional}@ = T;
    @\removed{using iterator   = \impdefnc;  // present only if \textbackslash{}tcode\{T\} is an object type other than an array}@
                                                @\removed{// of unknown bound; see\textasciitilde{}\ref{optional.ref.iterators}}@

    @\added{using iterator = \impdefnc; // see\textasciitilde{}\ref{optional.ref.iterators}}@

    // \ref{optional.ref.ctor}, constructors

    constexpr optional() noexcept = default;
    @\removed{constexpr optional(nullopt_t) noexcept : optional() \{\}}@
    @\removed{constexpr optional(const optional\& rhs) noexcept = default;}@

    @\removed{template<class Arg>}@
      @\removed{constexpr explicit optional(in_place_t, Arg\&\& arg);}@
    @\added{constexpr optional(nullopt_t) noexcept;}@

    @\added{constexpr optional(const optional\&) noexcept = default;}@

    @\added{template<class Arg> constexpr explicit optional(in_place_t, Arg\&\&);}@

    template<class U>
      @\removed{constexpr explicit(\seebelow) optional(U\&\& u) noexcept(\seebelow);}@
    @\added{constexpr explicit(!is_convertible_v<U, T\&>)}@
        @\added{optional(U\&\&) noexcept(is_nothrow_constructible_v<T\&, U>);}@

    template<class U>
    constexpr explicit(@\removed{\seebelow) optional(optional<U>\& rhs) noexcept(\seebelow);}@@\added{!is_convertible_v<U\&, T\&>)}@
        @\added{optional(optional<U>\&) noexcept(is_nothrow_constructible_v<T\&, U\&>);}@

    template<class U>
    constexpr explicit(@\removed{\seebelow) optional(const optional<U>\& rhs) noexcept(\seebelow);}@@\added{!is_convertible_v<const U\&, T\&>)}@
        @\added{optional(const optional<U>\&) noexcept(is_nothrow_constructible_v<T\&, const U\&>);}@

    template<class U>
    constexpr explicit(@\removed{\seebelow) optional(optional<U>\&\& rhs) noexcept(\seebelow);}@@\added{!is_convertible_v<U, T\&>)}@
        @\added{optional(optional<U>\&\&) noexcept(noexcept(is_nothrow_constructible_v<T\&, U>));}@

    template<class U>
    constexpr explicit(@\removed{\seebelow) optional(const optional<U>\&\& rhs) noexcept(\seebelow);}@@\added{!is_convertible_v<const U, T\&>)}@
        @\added{optional(const optional<U>\&\&) noexcept(}@
            @\added{noexcept(is_nothrow_constructible_v<T\&, const U>));}@

    constexpr ~optional() = default;

    // \ref{optional.ref.assign}, assignment
    constexpr optional& operator=(nullopt_t) noexcept;
    @\removed{constexpr optional\& operator=(const optional\& rhs) noexcept = default;}@

    @\removed{template<class U> constexpr T\& emplace(U\&\& u) noexcept(\seebelow);}@
    @\added{constexpr optional\& operator=(const optional\&) noexcept = default;}@

    @\added{template<class U>}@
    @\added{constexpr T\& emplace(U\&\&) noexcept(is_nothrow_constructible_v<T\&, U>);}@

    // \ref{optional.ref.swap}, swap
    constexpr void swap(optional& @\removed{rhs}@) noexcept;

    // \ref{optional.ref.iterators}, iterator support
    constexpr @\removed{auto}@@\added{iterator}@ begin() const noexcept;
    @\removed{constexpr auto end() const noexcept;}@

    @\added{constexpr iterator end() const noexcept;}@

    // \ref{optional.ref.observe}, observers
    constexpr T* operator->() const noexcept;

    constexpr T& operator*() const noexcept;

    constexpr explicit operator bool() const noexcept;
    @\removed{constexpr bool     has_value() const noexcept;}@
    @\removed{constexpr T\&       value() const;                           // freestanding-deleted}@
    @\removed{template<class U = remove_cv_t<T>>}@
      @\removed{constexpr remove_cv_t<T> value_or(U\&\& u) const;}@

    @\added{constexpr bool has_value() const noexcept;}@

    @\added{constexpr T\& value() const;}@

    @\added{template<class U = remove_cv_t<T>> constexpr decay_t<T> value_or(U\&\&) const;}@

    // \ref{optional.ref.monadic}, monadic operations
    template<class F> constexpr auto and_then(F&& @\removed{f}@) const;
    @\removed{template<class F> constexpr optional<invoke_result_t<F, T\&>> transform(F\&\& f) const;}@
    @\removed{template<class F> constexpr optional or_else(F\&\& f) const;}@

    @\added{template<class F> constexpr optional<invoke_result_t<F, T\&>> transform(F\&\&) const;}@

    @\added{template<class F> constexpr optional or_else(F\&\&) const;}@

    // \ref{optional.ref.mod}, modifiers
    constexpr void reset() noexcept;

  private:
    T* @\exposidnc{val}@ = nullptr; // \expos

    // \ref{optional.ref.expos}, exposition only helper functions
    template<class U> constexpr void @\exposidnc{convert-ref-init-val}@(U&& @\removed{u}@); // \expos
  };
}
\end{codeblock}

\pnum
An object of type \tcode{optional<T\&>} \removed{\defnx{contains a value}{contains a value!\idxcode{optional.ref}}}\added{\defn{contains a value}} if and only if \tcode{\exposid{val} != nullptr} is \tcode{true}. When an \tcode{optional<T\&>} contains a value, the \removed{\defnx{contained value}{contained value!\idxcode{optional.ref}}}\added{\defn{contained value}} is a reference to \tcode{*\exposid{val}}.

\pnum
Each type \tcode{optional<T\&>} is a trivially copyable class\iref{class.prop}.

\rSec3[optional.ref.ctor]{Constructors}

\begin{addedblock}
\indexlibraryctor{optional}%
\begin{itemdecl}
constexpr optional(nullopt_t) noexcept;
\end{itemdecl}

\begin{itemdescr}
\pnum
\ensures
\tcode{*this} does not contain a value.
\end{itemdescr}
\end{addedblock}

\indexlibraryctor{optional}%
\begin{itemdecl}
template<class Arg>
  constexpr explicit optional(in_place_t, Arg&& arg);
\end{itemdecl}

\begin{itemdescr}
\pnum
\constraints
\begin{itemize}
\item \tcode{is_constructible_v<T\&, Arg>} is \tcode{true}, and
\item \tcode{reference_constructs_from_temporary_v<T\&, Arg>} is \tcode{false}.
\end{itemize}

\pnum
\effects
Equivalent to: \tcode{\exposid{convert-ref-init-val}(std::forward<Arg>(arg))}.

\pnum
\ensures
\tcode{*this} contains a value.
\end{itemdescr}

\indexlibraryctor{optional}%
\begin{itemdecl}
template<class U>
  constexpr explicit(!is_convertible_v<U, T&>) optional(U&& u) noexcept(is_nothrow_constructible_v<T&, U>);
\end{itemdecl}

\begin{itemdescr}
\pnum
\constraints
\begin{itemize}
\item \tcode{is_same_v<remove_cvref_t<U>, optional>} is \tcode{false},
\item \tcode{is_same_v<remove_cvref_t<U>, in_place_t>} is \tcode{false}, and
\item \tcode{is_constructible_v<T\&, U>} is \tcode{true}.
\end{itemize}

\pnum
\effects
Equivalent to: \tcode{\exposid{convert-ref-init-val}(std::forward<U>(u))}.

\pnum
\ensures
\tcode{*this} contains a value.

\pnum
\remarks
This constructor is defined as deleted if
\begin{codeblock}
reference_constructs_from_temporary_v<T&, U>
\end{codeblock}
is \tcode{true}.
\end{itemdescr}

\indexlibraryctor{optional}%
\begin{itemdecl}
template<class U>
  constexpr explicit(!is_convertible_v<U&, T&>) optional(optional<U>& rhs) noexcept(is_nothrow_constructible_v<T&, U&>);
\end{itemdecl}

\begin{itemdescr}
\pnum
\constraints
\begin{itemize}
\item \tcode{is_same_v<remove_cv_t<T>, optional<U>>} is \tcode{false},
\item \tcode{is_same_v<T\&, U>} is \tcode{false}, and
\item \tcode{is_constructible_v<T\&, U\&>} is \tcode{true}.
\end{itemize}

\pnum
\effects
Equivalent to:
\begin{codeblock}
if (rhs.has_value()) @\exposid{convert-ref-init-val}@(rhs.operator*());
\end{codeblock}

\pnum
\remarks
This constructor is defined as deleted if
\begin{codeblock}
reference_constructs_from_temporary_v<T&, U&>
\end{codeblock}
is \tcode{true}.
\end{itemdescr}

\indexlibraryctor{optional}%
\begin{itemdecl}
template<class U>
  constexpr explicit(!is_convertible_v<const U&, T&>)
      optional(const optional<U>& rhs) noexcept(is_nothrow_constructible_v<T&, const U&>);
\end{itemdecl}

\begin{itemdescr}
\pnum
\constraints
\begin{itemize}
\item \tcode{is_same_v<remove_cv_t<T>, optional<U>>} is \tcode{false},
\item \tcode{is_same_v<T\&, U>} is \tcode{false}, and
\item \tcode{is_constructible_v<T\&, const U\&>} is \tcode{true}.
\end{itemize}

\pnum
\effects
Equivalent to:
\begin{codeblock}
if (rhs.has_value()) @\exposid{convert-ref-init-val}@(rhs.operator*());
\end{codeblock}

\pnum
\remarks
This constructor is defined as deleted if
\begin{codeblock}
reference_constructs_from_temporary_v<T&, const U&>
\end{codeblock}
is \tcode{true}.
\end{itemdescr}

\indexlibraryctor{optional}%
\begin{itemdecl}
template<class U>
  constexpr explicit(!is_convertible_v<U, T&>)
      optional(optional<U>&& rhs) noexcept(@\added{noexcept(}@is_nothrow_constructible_v<T&, U>)@\added{)}@;
\end{itemdecl}

\begin{itemdescr}
\pnum
\constraints
\begin{itemize}
\item \tcode{is_same_v<remove_cv_t<T>, optional<U>>} is \tcode{false},
\item \tcode{is_same_v<T\&, U>} is \tcode{false}, and
\item \tcode{is_constructible_v<T\&, U>} is \tcode{true}.
\end{itemize}

\pnum
\effects
Equivalent to:
\begin{codeblock}
if (rhs.has_value()) @\exposid{convert-ref-init-val}@(std::move(rhs).operator*());
\end{codeblock}

\pnum
\remarks
This constructor is defined as deleted if
\begin{codeblock}
reference_constructs_from_temporary_v<T&, U>
\end{codeblock}
is \tcode{true}.
\end{itemdescr}

\indexlibraryctor{optional}%
\begin{itemdecl}
template<class U>
  constexpr explicit(!is_convertible_v<const U, T&>)
      optional(const optional<U>&& rhs) noexcept(@\added{noexcept(}@is_nothrow_constructible_v<T&, const U>)@\added{)}@;
\end{itemdecl}

\begin{itemdescr}
\pnum
\constraints
\begin{itemize}
\item \tcode{is_same_v<remove_cv_t<T>, optional<U>>} is \tcode{false},
\item \tcode{is_same_v<T\&, U>} is \tcode{false}, and
\item \tcode{is_constructible_v<T\&, const U>} is \tcode{true}.
\end{itemize}

\pnum
\effects
Equivalent to:
\begin{codeblock}
if (rhs.has_value()) @\exposid{convert-ref-init-val}@(std::move(rhs).operator*());
\end{codeblock}

\pnum
\remarks
This constructor is defined as deleted if
\begin{codeblock}
reference_constructs_from_temporary_v<T&, const U>
\end{codeblock}
is \tcode{true}.
\end{itemdescr}

\rSec3[optional.ref.assign]{Assignment}

\indexlibrarymember{operator=}{optional}%
\begin{itemdecl}
constexpr optional& operator=(nullopt_t) noexcept;
\end{itemdecl}

\begin{itemdescr}
\pnum
\effects
Assigns \tcode{nullptr} to \exposid{val}.

\pnum
\ensures
\tcode{*this} does not contain a value.

\pnum
\returns
\tcode{*this}.
\end{itemdescr}

\indexlibrarymember{emplace}{optional}%
\begin{itemdecl}
template<class U>
  constexpr T& emplace(U&& u) noexcept(is_nothrow_constructible_v<T&, U>);
\end{itemdecl}

\begin{itemdescr}
\pnum
\constraints
\begin{itemize}
\item \tcode{is_constructible_v<T\&, U>} is \tcode{true}, and
\item \tcode{reference_constructs_from_temporary_v<T\&, U>} is \tcode{false}.
\end{itemize}

\pnum
\effects
Equivalent to: \tcode{\exposid{convert-ref-init-val}(std::forward<U>(u))}.

\pnum
\returns
\tcode{*\exposid{val}}.
\end{itemdescr}

\rSec3[optional.ref.swap]{Swap}

\indexlibrarymember{swap}{optional}%
\begin{itemdecl}
constexpr void swap(optional& rhs) noexcept;
\end{itemdecl}

\begin{itemdescr}
\pnum
\effects
Equivalent to: \tcode{std::swap(\exposid{val}, rhs.\exposid{val})}.
\end{itemdescr}

\rSec3[optional.ref.iterators]{Iterator support}

\begin{removedblock}
\begin{itemdecl}
using iterator = @\impdefnc@;        // present only if \tcode{T} is an object type other than an array
// of unknown bound
\end{itemdecl}

\begin{itemdescr}
\pnum
This type models \libconcept{contiguous_iterator}\iref{iterator.concept.contiguous}, meets the \oldconcept{RandomAccessIterator} requirements\iref{random.access.iterators}, and meets the requirements for constexpr iterators\iref{iterator.requirements.general}, with value type \tcode{remove_cv_t<T>}. The reference type is \tcode{T\&} for \tcode{iterator}.

\pnum
All requirements on container iterators\iref{container.reqmts} apply to \tcode{optional::iterator}.
\end{itemdescr}
\end{removedblock}

\begin{addedblock}
\indexlibrarymember{iterator}{optional}%
\begin{itemdecl}
using iterator = @\impdef@;
\end{itemdecl}

\begin{itemdescr}
\pnum
This type models \libconcept{contiguous_iterator}\iref{iterator.concept.contiguous}, meets the \oldconcept{RandomAccessIterator} requirements\iref{random.access.iterators}, and meets the requirements for constexpr iterators\iref{iterator.requirements.general}, with value type \tcode{remove_cv_t<T>}. The reference type is \tcode{T\&} for \tcode{iterator}.

\pnum
All requirements on container iterators\iref{container.reqmts} apply to \tcode{optional::iterator}.
\end{itemdescr}
\end{addedblock}

\indexlibrarymember{begin}{optional}%
\begin{itemdecl}
constexpr @\removed{auto}@@\added{iterator}@ begin() const noexcept;
\end{itemdecl}

\begin{itemdescr}
\begin{removedblock}
\pnum
\constraints
\tcode{T} is an object type other than an array of unknown bound.
\end{removedblock}

\pnum
\returns
An object \tcode{i} of type \tcode{iterator}, such that \tcode{i} is an iterator referring to \tcode{*\exposid{val}} if \tcode{has_value()} is \tcode{true}, and a past-the-end iterator value otherwise.
\end{itemdescr}

\indexlibrarymember{end}{optional}%
\begin{itemdecl}
constexpr @\removed{auto}@@\added{iterator}@ end() const noexcept;
\end{itemdecl}

\begin{itemdescr}
\begin{removedblock}
\pnum
\constraints
\tcode{T} is an object type other than an array of unknown bound.
\end{removedblock}

\pnum
\returns
\tcode{begin() + has_value()}.
\end{itemdescr}

\rSec3[optional.ref.observe]{Observers}

\indexlibrarymember{operator->}{optional}%
\begin{itemdecl}
constexpr T* operator->() const noexcept;
\end{itemdecl}

\begin{itemdescr}
\pnum
\hardexpects
\tcode{has_value()} is \tcode{true}.

\pnum
\returns
\exposid{val}.
\end{itemdescr}

\indexlibrarymember{operator*}{optional}%
\begin{itemdecl}
constexpr T& operator*() const noexcept;
\end{itemdecl}

\begin{itemdescr}
\pnum
\hardexpects
\tcode{has_value()} is \tcode{true}.

\pnum
\returns
\tcode{*\exposid{val}}.
\end{itemdescr}

\indexlibrarymember{operator bool}{optional}%
\begin{itemdecl}
constexpr explicit operator bool() const noexcept;
\end{itemdecl}

\begin{itemdescr}
\pnum
\returns
\tcode{\exposid{val} != nullptr}.
\end{itemdescr}

\indexlibrarymember{has_value}{optional}%
\begin{itemdecl}
constexpr bool has_value() const noexcept;
\end{itemdecl}

\begin{itemdescr}
\pnum
\returns
\tcode{\exposid{val} != nullptr}.
\end{itemdescr}

\indexlibrarymember{value}{optional}%
\begin{itemdecl}
constexpr T& value() const;
\end{itemdecl}

\begin{itemdescr}
\pnum
\effects
Equivalent to:
\begin{codeblock}
@\removed{return has_value() ? *\exposid{val} : throw bad_optional_access();}@
@\added{if (has_value()) \{}@
  @\added{return *\exposid{val};}@
@\added{\}}@
@\added{throw bad_optional_access();}@
\end{codeblock}
\end{itemdescr}

\indexlibrarymember{value_or}{optional}%
\begin{itemdecl}
template<class U = remove_cv_t<T>>
  constexpr @\removed{remove_cv_t}@@\added{decay_t}@<T> value_or(U&& u) const;
\end{itemdecl}

\begin{itemdescr}
\pnum
\constraints
\tcode{T} is a non-array object type.

\begin{removedblock}
\pnum
Let \tcode{X} be \tcode{remove_cv_t<T>}.
\end{removedblock}

\pnum
\mandates
\tcode{\removed{is_constructible_v<X, T\&> \&\& is_convertible_v<U, X>}}\tcode{\added{is_convertible_v<T\&, X>}} is \tcode{true} \added{and }\tcode{\added{is_convertible_v<U, X>}}\added{ is }\tcode{\added{true}}.

\pnum
\effects
Equivalent to:
\begin{codeblock}
@\removed{return has_value() ? *\exposid{val} : static_cast<X>(std::forward<U>(u));}@
@\added{using X = remove_cv_t<T>;}@

@\added{if (has_value()) \{}@
  @\added{return *\exposid{val};}@
@\added{\}}@
@\added{return std::forward<U>(u);}@
\end{codeblock}

\pnum
\remarks
The return type is unspecified if \tcode{T} is an array type or a non-object type.
\begin{note}
This is to avoid the declaration being ill-formed.
\end{note}
\end{itemdescr}

\rSec3[optional.ref.monadic]{Monadic operations}

\indexlibrarymember{and_then}{optional}%
\begin{itemdecl}
template<class F>
  constexpr auto and_then(F&& f) const;
\end{itemdecl}

\begin{itemdescr}
\pnum
Let \tcode{U} be \tcode{invoke_result_t<F, T\&>}.

\pnum
\mandates
\tcode{\removed{remove_cvref_t<U>}}\tcode{\added{U}} is a specialization of \tcode{optional}.

\pnum
\effects
Equivalent to:
\begin{codeblock}
if (has_value()) {
  return invoke(std::forward<F>(f), *@\exposid{val}@);
} else {
  return remove_cvref_t<U>();
}
\end{codeblock}
\end{itemdescr}

\indexlibrarymember{transform}{optional}%
\begin{itemdecl}
template<class F>
  constexpr optional<@\removed{remove_cv_t<}@invoke_result_t<F, T&>>@\removed{>}@ transform(F&& f) const;
\end{itemdecl}

\begin{itemdescr}
\pnum
Let \tcode{U} be \tcode{\removed{remove_cv_t<invoke_result_t<F, T\&>>}}\tcode{\added{invoke_result_t<F, T\&>}}.

\pnum
\mandates
\tcode{U} is a valid contained type for \tcode{optional}. The declaration
\begin{codeblock}
U u(invoke(std::forward<F>(f), *@\exposid{val}@));
\end{codeblock}
is well-formed for some invented variable \tcode{u}.
\begin{note}
There is no requirement that \tcode{U} is movable\iref{dcl.init.general}.
\end{note}

\pnum
\returns
If \tcode{*this} contains a value, an \tcode{optional<U>} object whose contained value is direct-non-list-initialized with \tcode{invoke(std::forward<F>(f), *\exposid{val})}; otherwise, \tcode{optional<U>()}.
\end{itemdescr}

\indexlibrarymember{or_else}{optional}%
\begin{itemdecl}
template<class F>
  constexpr optional or_else(F&& f) const;
\end{itemdecl}

\begin{itemdescr}
\pnum
\constraints
\tcode{F} models \libconcept{invocable}.

\pnum
\mandates
\tcode{is_same_v<remove_cvref_t<invoke_result_t<F>>, optional>} is \tcode{true}.

\pnum
\effects
Equivalent to:
\begin{codeblock}
if (has_value()) {
  return *@\removed{\exposid{val}}@@\added{this}@;
} else {
  return std::forward<F>(f)();
}
\end{codeblock}
\end{itemdescr}

\rSec3[optional.ref.mod]{Modifiers}

\indexlibrarymember{reset}{optional}%
\begin{itemdecl}
constexpr void reset() noexcept;
\end{itemdecl}

\begin{itemdescr}
\pnum
\effects
Assigns \tcode{nullptr} to \exposid{val}.

\pnum
\ensures
\tcode{*this} does not contain a value.
\end{itemdescr}

\rSec3[optional.ref.expos]{Exposition only helper functions}

\indexlibrarymember{convert_ref_init_val}{optional}%
\begin{itemdecl}
template<class U>
  constexpr void @\exposid{convert-ref-init-val}@(U&& u);  @\removed{// exposition only}@
\end{itemdecl}

\begin{itemdescr}
\pnum
\effects
Creates a variable \tcode{r} as if by \tcode{T\& r(std::forward<U>(u));} and then initializes \exposid{val} with \tcode{addressof(r)}.
\end{itemdescr}

\rSec2[optional.nullopt]{No-value state indicator}

\begin{removedblock}
\indexlibraryglobal{nullopt_t}%
\indexlibraryglobal{nullopt}%
\begin{itemdecl}
struct nullopt_t{@\seebelow@};
inline constexpr nullopt_t nullopt(@\unspec@);
\end{itemdecl}
\end{removedblock}

\begin{removedblock}
\pnum
The struct \tcode{nullopt_t} is an empty class type used as a unique type to indicate the state of not containing a value for \tcode{optional} objects. In particular, \tcode{optional<T>} has a constructor with \tcode{nullopt_t} as a single argument; this indicates that an optional object not containing a value shall be constructed.

\pnum
Type \tcode{nullopt_t} does not have a default constructor or an initializer-list constructor, and is not an aggregate. \tcode{nullopt_t} models \libconcept{copyable} and \tcode{\libconcept{three_way_comparable}<strong_ordering>}.
\end{removedblock}

\begin{addedblock}
\indexlibraryglobal{nullopt}%
\begin{itemdecl}
struct nullopt_t{@\seebelow@};
inline constexpr nullopt_t nullopt(@\unspec@);
\end{itemdecl}

\begin{itemdescr}
\pnum
The struct \tcode{nullopt_t} is an empty class type used as a unique type to indicate the state of not containing a value for \tcode{optional} objects. In particular, \tcode{optional<T>} has a constructor with \tcode{nullopt_t} as a single argument; this indicates that an optional object not containing a value shall be constructed.

\pnum
Type \tcode{nullopt_t} does not have a default constructor or an initializer-list constructor, and is not an aggregate. \tcode{nullopt_t} models \libconcept{copyable} and \tcode{three_way_comparable<strong_ordering>}.
\end{itemdescr}
\end{addedblock}

\begin{removedblock}
\rSec2[optional.bad.access]{Class \tcode{bad_optional_access}}

\begin{codeblock}
namespace std {
  class bad_optional_access : public exception {
  public:
    // see \ref{exception} for the specification of the special member functions
    constexpr const char* what() const noexcept override;
  };
}
\end{codeblock}

\pnum
The class \tcode{bad_optional_access} defines the type of objects thrown as exceptions to report the situation where an attempt is made to access the value of an optional object that does not contain a value.

\indexlibrarymember{what}{bad_optional_access}%
\begin{itemdecl}
constexpr const char* what() const noexcept override;
\end{itemdecl}

\begin{itemdescr}
\pnum
\returns
An \impldef{return value of \tcode{bad_optional_access::what}} \ntbs{}, which during constant evaluation is encoded with the ordinary literal encoding\iref{lex.ccon}.
\end{itemdescr}
\end{removedblock}

\rSec2[optional.relops]{Relational operators}

\indexlibraryglobal{operator==}%
\begin{itemdecl}
template<class T, class U>
  constexpr bool operator==(const optional<T>& x, const optional<U>& y);
\end{itemdecl}

\begin{itemdescr}
\pnum
\constraints
The expression \tcode{*x == *y} is well-formed and its result is convertible to \tcode{bool}.
\begin{note}
\tcode{T} need not be \oldconcept{EqualityComparable}.
\end{note}

\pnum
\returns
If \tcode{x.has_value() != y.has_value()}, \tcode{false}; otherwise if \tcode{x.has_value() == false}, \tcode{true}; otherwise \tcode{*x == *y}.

\pnum
\remarks
Specializations of this function template for which \tcode{*x == *y} is a core constant expression are constexpr functions.
\end{itemdescr}

\begin{removedblock}
\indexlibrarymember{operator"!=}{optional}%
\begin{itemdecl}
template<class T, class U> constexpr bool operator!=(const optional<T>& x, const optional<U>& y);
\end{itemdecl}

\begin{itemdescr}
\pnum
\constraints
The expression \tcode{*x != *y} is well-formed and its result is convertible to \tcode{bool}.

\pnum
\returns
If \tcode{x.has_value() != y.has_value()}, \tcode{true}; otherwise, if \tcode{x.has_value() == false}, \tcode{false}; otherwise \tcode{*x != *y}.

\pnum
\remarks
Specializations of this function template for which \tcode{*x != *y} is a core constant expression are constexpr functions.
\end{itemdescr}
\end{removedblock}

\begin{addedblock}
\indexlibraryglobal{operator!=}%
\begin{itemdecl}
template<class T, class U>
  constexpr bool operator!=(const optional<T>& x, const optional<U>& y);
\end{itemdecl}

\begin{itemdescr}
\pnum
\constraints
The expression \tcode{*x != *y} is well-formed and its result is convertible to \tcode{bool}.

\pnum
\returns
If \tcode{x.has_value() != y.has_value()}, \tcode{true}; otherwise, if \tcode{x.has_value() == false}, \tcode{false}; otherwise \tcode{*x != *y}.

\pnum
\remarks
Specializations of this function template for which \tcode{*x != *y} is a core constant expression are constexpr functions.
\end{itemdescr}
\end{addedblock}

\indexlibraryglobal{operator<}%
\begin{itemdecl}
template<class T, class U>
  constexpr bool operator<(const optional<T>& x, const optional<U>& y);
\end{itemdecl}

\begin{itemdescr}
\pnum
\constraints
\tcode{*x < *y} is well-formed and its result is convertible to \tcode{bool}.

\pnum
\returns
If \tcode{!y}, \tcode{false}; otherwise, if \tcode{!x}, \tcode{true}; otherwise \tcode{*x < *y}.

\pnum
\remarks
Specializations of this function template for which \tcode{*x < *y} is a core constant expression are constexpr functions.
\end{itemdescr}

\indexlibraryglobal{operator>}%
\begin{itemdecl}
template<class T, class U>
  constexpr bool operator>(const optional<T>& x, const optional<U>& y);
\end{itemdecl}

\begin{itemdescr}
\pnum
\constraints
The expression \tcode{*x > *y} is well-formed and its result is convertible to \tcode{bool}.

\pnum
\returns
If \tcode{!x}, \tcode{false}; otherwise, if \tcode{!y}, \tcode{true}; otherwise \tcode{*x > *y}.

\pnum
\remarks
Specializations of this function template for which \tcode{*x > *y} is a core constant expression are constexpr functions.
\end{itemdescr}

\indexlibraryglobal{operator<=}%
\begin{itemdecl}
template<class T, class U>
  constexpr bool operator<=(const optional<T>& x, const optional<U>& y);
\end{itemdecl}

\begin{itemdescr}
\pnum
\constraints
The expression \tcode{*x <= *y} is well-formed and its result is convertible to \tcode{bool}.

\pnum
\returns
If \tcode{!x}, \tcode{true}; otherwise, if \tcode{!y}, \tcode{false}; otherwise \tcode{*x <= *y}.

\pnum
\remarks
Specializations of this function template for which \tcode{*x <= *y} is a core constant expression are constexpr functions.
\end{itemdescr}

\indexlibraryglobal{operator>=}%
\begin{itemdecl}
template<class T, class U>
  constexpr bool operator>=(const optional<T>& x, const optional<U>& y);
\end{itemdecl}

\begin{itemdescr}
\pnum
\constraints
The expression \tcode{*x >= *y} is well-formed and its result is convertible to \tcode{bool}.

\pnum
\returns
If \tcode{!y}, \tcode{true}; otherwise, if \tcode{!x}, \tcode{false}; otherwise \tcode{*x >= *y}.

\pnum
\remarks
Specializations of this function template for which \tcode{*x >= *y} is a core constant expression are constexpr functions.
\end{itemdescr}

\indexlibraryglobal{operator<=>}%
\begin{itemdecl}
template<class T, @\removed{@\textbackslash{}libconcept\{three_way_comparable_with\}@}@@\added{three_way_comparable_with}@<T> U>
  constexpr compare_three_way_result_t<T, U> operator<=>(const optional<T>& x, const optional<U>& y);
\end{itemdecl}

\begin{itemdescr}
\pnum
\returns
If \tcode{x \&\& y}, \tcode{*x <=> *y}; otherwise \tcode{x.has_value() <=> y.has_value()}.

\pnum
\remarks
Specializations of this function template for which \tcode{*x <=> *y} is a core constant expression are constexpr functions.
\end{itemdescr}

\rSec2[optional.nullops]{Comparison with \tcode{nullopt}}

\indexlibraryglobal{operator==}%
\begin{itemdecl}
template<class T>
  constexpr bool operator==(const optional<T>& x, nullopt_t) noexcept;
\end{itemdecl}

\begin{itemdescr}
\pnum
\returns
\tcode{!x}.
\end{itemdescr}

\indexlibraryglobal{operator<=>}%
\begin{itemdecl}
template<class T>
  constexpr strong_ordering operator<=>(const optional<T>& x, nullopt_t) noexcept;
\end{itemdecl}

\begin{itemdescr}
\pnum
\returns
\tcode{x.has_value() <=> false}.
\end{itemdescr}

\rSec2[optional.comp.with.t]{Comparison with \tcode{T}}

\indexlibraryglobal{operator==}%
\begin{itemdecl}
template<class T, class U>
  constexpr bool operator==(const optional<T>& x, const U& v);
\end{itemdecl}

\begin{itemdescr}
\pnum
\constraints
\tcode{U} is not a specialization of \tcode{optional}. The expression \tcode{*x == v} is well-formed and its result is convertible to \tcode{bool}.
\begin{note}
\tcode{T} need not be \oldconcept{EqualityComparable}.
\end{note}

\pnum
\effects
Equivalent to:
\begin{codeblock}
if (x.has_value())
  return *x == v;
return false;
\end{codeblock}
\end{itemdescr}

\indexlibraryglobal{operator==}%
\begin{itemdecl}
template<class T, class U>
  constexpr bool operator==(const T& v, const optional<U>& x);
\end{itemdecl}

\begin{itemdescr}
\pnum
\constraints
\tcode{T} is not a specialization of \tcode{optional}. The expression \tcode{v == *x} is well-formed and its result is convertible to \tcode{bool}.

\pnum
\effects
Equivalent to:
\begin{codeblock}
if (x.has_value())
  return v == *x;
return false;
\end{codeblock}
\end{itemdescr}

\begin{removedblock}
\indexlibrarymember{operator"!=}{optional}%
\begin{itemdecl}
template<class T, class U> constexpr bool operator!=(const optional<T>& x, const U& v);
\end{itemdecl}

\begin{itemdescr}
\pnum
\constraints
\tcode{U} is not a specialization of \tcode{optional}. The expression \tcode{*x != v} is well-formed and its result is convertible to \tcode{bool}.

\pnum
\effects
Equivalent to:
\begin{codeblock}
if (x.has_value())
  return *x != v;
return true;
\end{codeblock}
\end{itemdescr}
\end{removedblock}

\begin{removedblock}
\indexlibrarymember{operator"!=}{optional}%
\begin{itemdecl}
template<class T, class U> constexpr bool operator!=(const T& v, const optional<U>& x);
\end{itemdecl}

\begin{itemdescr}
\pnum
\constraints
\tcode{T} is not a specialization of \tcode{optional}. The expression \tcode{v != *x} is well-formed and its result is convertible to \tcode{bool}.

\pnum
\effects
Equivalent to:
\begin{codeblock}
if (x.has_value())
  return v != *x;
return true;
\end{codeblock}
\end{itemdescr}
\end{removedblock}

\begin{addedblock}
\indexlibraryglobal{operator!=}%
\begin{itemdecl}
template<class T, class U>
  constexpr bool operator!=(const optional<T>& x, const U& v);
\end{itemdecl}

\begin{itemdescr}
\pnum
\constraints
\tcode{U} is not a specialization of \tcode{optional}. The expression \tcode{*x != v} is well-formed and its result is convertible to \tcode{bool}.

\pnum
\effects
Equivalent to:
\begin{codeblock}
if (x.has_value())
  return *x != v;
return true;
\end{codeblock}
\end{itemdescr}
\end{addedblock}

\begin{addedblock}
\indexlibraryglobal{operator!=}%
\begin{itemdecl}
template<class T, class U>
  constexpr bool operator!=(const T& v, const optional<U>& x);
\end{itemdecl}

\begin{itemdescr}
\pnum
\constraints
\tcode{T} is not a specialization of \tcode{optional}. The expression \tcode{v != *x} is well-formed and its result is convertible to \tcode{bool}.

\pnum
\effects
Equivalent to:
\begin{codeblock}
if (x.has_value())
  return v != *x;
return true;
\end{codeblock}
\end{itemdescr}
\end{addedblock}

\indexlibraryglobal{operator<}%
\begin{itemdecl}
template<class T, class U>
  constexpr bool operator<(const optional<T>& x, const U& v);
\end{itemdecl}

\begin{itemdescr}
\pnum
\constraints
\tcode{U} is not a specialization of \tcode{optional}. The expression \tcode{*x < v} is well-formed and its result is convertible to \tcode{bool}.

\pnum
\effects
Equivalent to:
\begin{codeblock}
if (x.has_value())
  return *x < v;
return true;
\end{codeblock}
\end{itemdescr}

\indexlibraryglobal{operator<}%
\begin{itemdecl}
template<class T, class U>
  constexpr bool operator<(const T& v, const optional<U>& x);
\end{itemdecl}

\begin{itemdescr}
\pnum
\constraints
\tcode{T} is not a specialization of \tcode{optional}. The expression \tcode{v < *x} is well-formed and its result is convertible to \tcode{bool}.

\pnum
\effects
Equivalent to:
\begin{codeblock}
if (x.has_value())
  return v < *x;
return false;
\end{codeblock}
\end{itemdescr}

\indexlibraryglobal{operator>}%
\begin{itemdecl}
template<class T, class U>
  constexpr bool operator>(const optional<T>& x, const U& v);
\end{itemdecl}

\begin{itemdescr}
\pnum
\constraints
\tcode{U} is not a specialization of \tcode{optional}. The expression \tcode{*x > v} is well-formed and its result is convertible to \tcode{bool}.

\pnum
\effects
Equivalent to:
\begin{codeblock}
if (x.has_value())
  return *x > v;
return false;
\end{codeblock}
\end{itemdescr}

\indexlibraryglobal{operator>}%
\begin{itemdecl}
template<class T, class U>
  constexpr bool operator>(const T& v, const optional<U>& x);
\end{itemdecl}

\begin{itemdescr}
\pnum
\constraints
\tcode{T} is not a specialization of \tcode{optional}. The expression \tcode{v > *x} is well-formed and its result is convertible to \tcode{bool}.

\pnum
\effects
Equivalent to:
\begin{codeblock}
if (x.has_value())
  return v > *x;
return true;
\end{codeblock}
\end{itemdescr}

\indexlibraryglobal{operator<=}%
\begin{itemdecl}
template<class T, class U>
  constexpr bool operator<=(const optional<T>& x, const U& v);
\end{itemdecl}

\begin{itemdescr}
\pnum
\constraints
\tcode{U} is not a specialization of \tcode{optional}. The expression \tcode{*x <= v} is well-formed and its result is convertible to \tcode{bool}.

\pnum
\effects
Equivalent to:
\begin{codeblock}
if (x.has_value())
  return *x <= v;
return true;
\end{codeblock}
\end{itemdescr}

\indexlibraryglobal{operator<=}%
\begin{itemdecl}
template<class T, class U>
  constexpr bool operator<=(const T& v, const optional<U>& x);
\end{itemdecl}

\begin{itemdescr}
\pnum
\constraints
\tcode{T} is not a specialization of \tcode{optional}. The expression \tcode{v <= *x} is well-formed and its result is convertible to \tcode{bool}.

\pnum
\effects
Equivalent to:
\begin{codeblock}
if (x.has_value())
  return v <= *x;
return false;
\end{codeblock}
\end{itemdescr}

\indexlibraryglobal{operator>=}%
\begin{itemdecl}
template<class T, class U>
  constexpr bool operator>=(const optional<T>& x, const U& v);
\end{itemdecl}

\begin{itemdescr}
\pnum
\constraints
\tcode{U} is not a specialization of \tcode{optional}. The expression \tcode{*x >= v} is well-formed and its result is convertible to \tcode{bool}.

\pnum
\effects
Equivalent to:
\begin{codeblock}
if (x.has_value())
  return *x >= v;
return false;
\end{codeblock}
\end{itemdescr}

\indexlibraryglobal{operator>=}%
\begin{itemdecl}
template<class T, class U>
  constexpr bool operator>=(const T& v, const optional<U>& x);
\end{itemdecl}

\begin{itemdescr}
\pnum
\constraints
\tcode{T} is not a specialization of \tcode{optional}. The expression \tcode{v >= *x} is well-formed and its result is convertible to \tcode{bool}.

\pnum
\effects
Equivalent to:
\begin{codeblock}
if (x.has_value())
  return v >= *x;
return true;
\end{codeblock}
\end{itemdescr}

\indexlibraryglobal{operator<=>}%
\begin{itemdecl}
template<class T, class U>
  requires(!@\removed{@\textbackslash{}exposconcept\{is-derived-from-optional\}@}@@\added{\exposid{is-derived-from-optional}}@<U>) && @\removed{@\textbackslash{}libconcept\{three_way_comparable_with\}@}@@\added{three_way_comparable_with}@<T, U>
  @\removed{constexpr compare_three_way_result_t<T, U>}@
    @\removed{operator<=>(const optional<T>\& x, const U\& v);}@
  @\added{constexpr compare_three_way_result_t<T, U> operator<=>(const optional<T>\& x, const U\& v);}@
\end{itemdecl}

\begin{itemdescr}
\pnum
\effects
Equivalent to: \tcode{return x.has_value() ? *x <=> v : strong_ordering::less;}
\end{itemdescr}

\rSec2[optional.specalg]{Specialized algorithms}

\indexlibraryglobal{swap}%
\begin{itemdecl}
template<class T>
  constexpr void swap(optional<T>& x, optional<T>& y) noexcept(noexcept(x.swap(y)));
\end{itemdecl}

\begin{itemdescr}
\begin{removedblock}
\pnum
\constraints
\begin{codeblock}
is_reference_v<T> || (is_move_constructible_v<T> && is_swappable_v<T>)
\end{codeblock}
\end{removedblock}
\begin{removedblock}
is \tcode{true}.
\end{removedblock}

\begin{addedblock}
\pnum
\tcode{is_move_constructible_v<T>} is \tcode{true} and \tcode{is_swappable_v<T>} is \tcode{true}.
\end{addedblock}

\pnum
\effects
Calls \tcode{x.swap(y)}.
\end{itemdescr}

\indexlibraryglobal{make_optional}%
\begin{itemdecl}
@\removed{template<class T> constexpr optional<decay_t<T>> make_optional(T\&\& v);}@
@\added{template<int = 0, class T>}@
  @\added{constexpr optional<decay_t<T>> make_optional(T\&\& v) noexcept(is_nothrow_constructible_v<optional<decay_t<T>>, T>);}@
\end{itemdecl}

\begin{itemdescr}
\pnum
\constraints
The call to \tcode{make_optional} does not use an explicit \removed{\grammarterm{template-argument-list}}\added{\grammarterm{template-argument-list}} that begins with a type \removed{\grammarterm{template-argument}}\added{\grammarterm{template-argument}, and }\tcode{\added{is_constructible_v<decay_t<T>, T>}}\added{ is }\tcode{\added{true}}.

\begin{removedblock}
\pnum
\effects
Equivalent to: \tcode{return optional<decay_t<T>>(std::forward<T>(v));}
\end{removedblock}

\begin{addedblock}
\pnum
Equivalent to:
\begin{codeblock}
return optional<decay_t<T>>{std::forward<T>(v)};
\end{codeblock}
\end{addedblock}
\end{itemdescr}

\indexlibraryglobal{make_optional}%
\begin{itemdecl}
template<class T, class... Args>
  constexpr optional<T> make_optional(Args&&... args) @\added{noexcept(is_nothrow_constructible_v<T, Args...>)}@;
\end{itemdecl}

\begin{itemdescr}
\begin{addedblock}
\pnum
\constraints
\tcode{is_constructible_v<T, Args...>} is \tcode{true}.
\end{addedblock}

\begin{addedblock}
\pnum
\mandates
\tcode{is_reference_v<T>} is \tcode{false}.
\end{addedblock}

\begin{removedblock}
\pnum
\effects
Equivalent to: \tcode{return optional<T>(in_place, std::forward<Args>(args)...);}
\end{removedblock}

\begin{addedblock}
\pnum
Equivalent to:
\begin{codeblock}
return optional<T>{in_place, std::forward<Args>(args)...};
\end{codeblock}
\end{addedblock}
\end{itemdescr}

\indexlibraryglobal{make_optional}%
\begin{itemdecl}
template<class T, class U, class... Args>
  constexpr optional<T> make_optional(initializer_list<U> il, Args&&... args) @\added{noexcept(}@
      @\added{is_nothrow_constructible_v<T, initializer_list<U>\&, Args...>)}@;
\end{itemdecl}

\begin{itemdescr}
\begin{addedblock}
\pnum
\constraints
\tcode{is_constructible_v<T, initializer_list<U>\&, Args...>} is \tcode{true}.
\end{addedblock}

\begin{addedblock}
\pnum
\mandates
\tcode{is_reference_v<U>} is \tcode{false}.
\end{addedblock}

\begin{removedblock}
\pnum
\effects
Equivalent to: \tcode{return optional<T>(in_place, il, std::forward<Args>(args)...);}
\end{removedblock}

\begin{addedblock}
\pnum
Equivalent to:
\begin{codeblock}
return optional<T>{in_place, il, std::forward<Args>(args)...};
\end{codeblock}
\end{addedblock}
\end{itemdescr}

\rSec2[optional.hash]{Hash support}

\begin{itemdecl}
template<class T> struct hash<optional<T>>;
\end{itemdecl}

\begin{itemdescr}
\pnum
The specialization \tcode{hash<optional<T>>} is enabled\iref{unord.hash} if and only if \tcode{hash<remove_const_t<T>>} is enabled. When enabled, for an object \tcode{o} of type \tcode{optional<T>}, if \tcode{o.has_value() == true}, then \tcode{hash<optional<T>>()(o)} evaluates to the same value as \tcode{hash<remove_const_t<T>>()(*o)}; otherwise it evaluates to an unspecified value. The member functions are not guaranteed to be \removed{\keyword{noexcept}}\tcode{\added{noexcept}}.
\end{itemdescr}
